\documentclass[12pt]{article}
\usepackage[utf8]{inputenc}
\usepackage[russian]{babel}
\usepackage{graphicx}
\usepackage{array}
\usepackage[a4paper, left=1cm, right=1cm, top=1.27cm, bottom=1.27cm]{geometry}
\usepackage{setspace}
\usepackage{tabularx}
\usepackage{makecell}
\usepackage{longtable}
\usepackage{multirow}
\usepackage{caption}
\onehalfspacing


\begin{document}


\begin{table}[h]
    \centering
    \renewcommand{\arraystretch}{1.5} % Увеличиваем межстрочный интервал
    \begin{tabular}{|m{11cm}|m{6cm}|} % Первая колонка шире, вторая для лого
        \hline
        \begin{center}
            \textbf{Университет ИТМО}\\
            \textbf{Физико-технический мегафакультет}\\
            \textbf{Физический факультет}
        \end{center} 
        & 
        \begin{center}
            
\includegraphics[height=1.3cm]{logo.jpg}
        \end{center} \\
        \hline
    \end{tabular}
\end{table}

\noindent\makebox[\linewidth]{\rule{\textwidth}{2pt}}

\vspace{10pt}

\begin{table}[h]
    \centering
    \renewcommand{\arraystretch}{1.7}
    \begin{tabular}{|p{4cm}p{4cm}|p{8cm}|}
        \hline
        Группа N3150 & & К работе допущен \underline{\hspace{4cm}} \\
        \hline
        Студенты & \hfill \multirow[0pt]{3}{*}{\makecell[c]{Пахомов В.Ю \\ Фармахини Ф.Р. \\ Петров Д.О}}  & Работа выполнена \underline{\hspace{4cm}} \\
        & & \\
        & & \\
        \hline
        \multicolumn{2}{|c|}{Преподаватель \underline{\hspace{5cm}}}  & Отчет принят \underline{\hspace{5cm}} \\
        \hline
    \end{tabular}
\end{table}

\begin{center}
    \fontsize{20pt}{21.5pt} \selectfont
    \textbf{Рабочий протокол и отчет по}\\
    \textbf{лабораторной работе № 2.01}
\end{center}

\noindent\makebox[\linewidth]{\rule{\textwidth}{1pt}}
\noindent\makebox[\linewidth]{\rule{\textwidth}{1pt}}

\begin{flushleft}
    1. Цель работы.\\
    \hspace{0.63cm}1) Экспериментальная проверка уравнения состояния идеального газа. \\
    \hspace{0.63cm}2) Определение температуры абсолютного нуля по шкале Цельсия. \\
\end{flushleft}

\begin{flushleft}
    2. Задачи, решаемые при выполнении работы.\\
    \hspace{0.63cm}1) Получить зависимости давления $p(V_c)$ при различных температурах $t$. \\
    \hspace{0.63cm}2) Построить графики зависимости $V_c(1/p)$ при различных
    температурах и $p(T)$ при разных значениях $V_c$. \\
\end{flushleft}

\begin{flushleft}
    3. Объект исследования. \\
\end{flushleft}

\begin{flushleft}
    4. Метод экспериментального исследования. \\
    \hspace{0.63cm}1) Изотермические процессы (закон Бойля-Мариотта) \\
    Измерение зависимости давления $p$ от объёма $V$ при фиксированных
    температурах (20–60°C). \\
    \hspace{0.63cm}2) Изохорные процессы (закон Шарля) \\
    Измерение зависимости давления $p$ от температуры $t$ при постоянных объёмах.

\end{flushleft}

\begin{flushleft}
    5. Рабочие формулы и исходные данные. \\
    $ p = p_0 + \frac{\Delta p_1 + \Delta p_2}{2} $ \\
    $ K = \frac{\Delta V_{c}}{\Delta (1/p)} = V_c \cdot p $ \\
    $ K = \frac{1}{D} \sum^N_{i=1}(X_i - \bar{X})Y_i, C = \bar{Y} - K\bar{X} $ \\
    $ X = 1/p, Y = V_c, \bar{X} $ и $ \bar{Y} $ - средние значения, $ D = \sum(X_i - \bar{X})^2 $ \\
    $ \Delta K = \sqrt{\frac{E}{D}}, E = \frac{1}{N-2}\sum^N_{i=1}(Y_i-KX_i-C)^2 $ \\
    $ t_* = -\frac{C}{A} $ \\
    $ \Delta t_* = t_* \sqrt{(\frac{\Delta A}{A})^2 + (\frac{\Delta C}{C})^2} $ \\
    $ A = \frac{1}{D} \sum_{i=1}^{N} (X_i - \bar{X}) Y_i, \quad C = \bar{Y} - A \bar{X} \label{eq:16} $ \\
    $ \bar{X} = \frac{1}{N} \sum_{i=1}^{N} X_i, \quad \bar{Y} = \frac{1}{N} \sum_{i=1}^{N} Y_i, \quad D = \sum_{i=1}^{N} (X_i - \bar{X})^2 $ \\
\end{flushleft}

\begin{flushleft}
    6. Измерительные приборы. \\
    \begin{table}[h]
        \centering
        \renewcommand{\arraystretch}{1.5}
        {\itshape
        \begin{tabular}{|c|c|c|c|c|}
            \hline
            № п/п & Наименование & Тип прибора & Используемый диапазон & Погрешность прибора\\
            \hline
            1 & ПКЦ-3 & Электронный & 0-99.99° & ±0.01° \\
            \hline
            2 & Цилиндр с поршнем & Механический & 50-120мл & ±0.1мл \\
            \hline
            3 & Барометр & Механический & 80-106 Па & ±1 Па \\
            \hline
        \end{tabular}
        }
    \end{table}
\end{flushleft}

\begin{flushleft}
    7. Схема установки. \\
    \includegraphics[height=8cm]{schema.png} \\
    \hspace{0.63cm}1)	цилиндр с поршнем \\
    \hspace{0.63cm}2)	опорная площадка цилиндра \\
    \hspace{0.63cm}3)	термостат \\
    \hspace{0.63cm}4)	щуп с датчиком температуры \\
    \hspace{0.63cm}5)	манометрический датчик \\
    \hspace{0.63cm}6)	стенд \\
    \hspace{0.63cm}7)	преобразователь сигналов \\
    \hspace{0.63cm}8)	измерительный прибор ПКЦ-3 \\
    \hspace{0.63cm}9)	кружка \\
    \hspace{0.63cm}10)	поддон \\
    \hspace{0.63cm}11)	лопатка \\
\end{flushleft}

\begin{flushleft}
    8. Результаты прямых измерений и их обработки. \\
    $p_0$ = 102.9 кПа \\
    \newpage

    \begin{center}
        \textbf{Таблица 1.1. Зависимость давления от объема при температуре $ t_1 = 21.7^\circ C $} \\
    \end{center}
    \begin{table}[h]
        \centering
        \begin{tabular}{|c|c|c|c|c|c|}
            \hline
            № п.п. & $ V_u $, мл & $ \Delta p_1 $, кПа & $ \Delta p_2 $, кПа & $ p $, кПа & $ 1/p $, кПа \\
            \hline
            1  & 50  & 53.4  & 53.4  & 156.3 & 0.00640 \\
            \hline
            2  & 60  & 28.0  & 28.6  & 131.5 & 0.00760 \\
            \hline
            3  & 70  & 10.3  & 12.8  & 114.0 & 0.00877 \\
            \hline
            4  & 80  & -3.7  & -0.1  & 99.6  & 0.01004 \\
            \hline
            5  & 90  & -13.9 & -12   & 89.0  & 0.01124 \\
            \hline
            6  & 100 & -22.2 & -21.5 & 80.2  & 0.01247 \\
            \hline
            7  & 110 & -28.7 & -28.7 & 73.2  & 0.01366 \\
            \hline
            8  & 120 & -35.2 & -35.2 & 66.7  & 0.01499 \\
            \hline
        \end{tabular}    
    \end{table} 
    
    \begin{center}
        \textbf{Таблица 1.2. Зависимость давления от объема при температуре $ t_2 = 30.5^\circ C $} \\
    \end{center}
    \begin{table}[h]
        \centering
        \begin{tabular}{|c|c|c|c|c|c|}
            \hline
            № п.п. & $ V_u $, мл & $ \Delta p_1 $, кПа & $ \Delta p_2 $, кПа & $ p $, кПа & $ 1/p $, кПа \\
            \hline
            1  & 50  & 53.5  & 53.5  & 156.4 & 0.00639 \\
            \hline
            2  & 60  & 29.5  & 33.3  & 132.9 & 0.00752 \\
            \hline
            3  & 70  & 12.0  & 15.5  & 115.2 & 0.00868 \\
            \hline
            4  & 80  & -0.8  & 0.8   & 102.9 & 0.00972 \\
            \hline
            5  & 90  & -11.5 & -10.5 & 91.4  & 0.01094 \\
            \hline
            6  & 100 & -19.1 & -20.1 & 82.4  & 0.01214 \\
            \hline
            7  & 110 & -27.2 & -27.3 & 74.7  & 0.01339 \\
            \hline
            8  & 120 & -33.5 & -33.5 & 68.4  & 0.01462 \\
            \hline
        \end{tabular}    
    \end{table}

    \begin{center}
        \textbf{Таблица 1.3. Зависимость давления от объема при температуре $ t_3 = 40.7^\circ C $} \\
    \end{center}
    \begin{table}[h]
        \centering
        \begin{tabular}{|c|c|c|c|c|c|}
            \hline
            № п.п. & $ V_u $, мл & $ \Delta p_1 $, кПа & $ \Delta p_2 $, кПа & $ p $, кПа & $ 1/p $, кПа \\
            \hline
            1  & 50  & 56.0  & 56.0  & 158.9 & 0.00629 \\
            \hline
            2  & 60  & 33.4  & 33.5  & 136.4 & 0.00733 \\
            \hline
            3  & 70  & 14.8  & 16.5  & 117.1 & 0.00854 \\
            \hline
            4  & 80  & 2.7   & 2.5   & 104.1 & 0.00961 \\
            \hline
            5  & 90  & -3.3  & -8.5  & 95.1  & 0.01052 \\
            \hline
            6  & 100 & -17.9 & -18.0 & 84.0  & 0.01190 \\
            \hline
            7  & 110 & -25.0 & -25.3 & 76.6  & 0.01305 \\
            \hline
            8  & 120 & -32.3 & -32.3 & 69.6  & 0.01437 \\
            \hline
        \end{tabular}    
    \end{table}

    \newpage
    \begin{center}
        \textbf{Таблица 1.4. Зависимость давления от объема при температуре $ t_4 = 49.8^\circ C $} \\
    \end{center}
    \begin{table}[h]
        \centering
        \begin{tabular}{|c|c|c|c|c|c|}
            \hline
            № п.п. & $ V_u $, мл & $ \Delta p_1 $, кПа & $ \Delta p_2 $, кПа & $ p $, кПа & $ 1/p $, кПа \\
            \hline
            1  & 50  & 62.7  & 62.7  & 165.6 & 0.00604 \\
            \hline
            2  & 60  & 38.1  & 36.0  & 139.0 & 0.00719 \\
            \hline
            3  & 70  & 18.6  & 19.3  & 120.2 & 0.00832 \\
            \hline
            4  & 80  & 4.0   & 4.4   & 106.6 & 0.00938 \\
            \hline
            5  & 90  & -7.1  & -6.8  & 95.9  & 0.01043 \\
            \hline
            6  & 100 & -16.5 & -14.3 & 87.1  & 0.01148 \\
            \hline
            7  & 110 & -23.0 & -22.7 & 79.2  & 0.01263 \\
            \hline
            8  & 120 & -30.2 & -30.2 & 71.7  & 0.01395 \\
            \hline
        \end{tabular}    
    \end{table}

    \begin{center}
        \textbf{Таблица 1.5. Зависимость давления от объема при температуре $ t_5 = 60^\circ C $} \\
    \end{center}
    \begin{table}[h]
        \centering
        \begin{tabular}{|c|c|c|c|c|c|}
            \hline
            № п.п. & $ V_u $, мл & $ \Delta p_1 $, кПа & $ \Delta p_2 $, кПа & $ p $, кПа & $ 1/p $, кПа \\
            \hline
            1  & 50  & 68    & 68    & 170.9 & 0.00585 \\
            \hline
            2  & 60  & 43    & 44.2  & 146.0 & 0.00685 \\
            \hline
            3  & 70  & 24    & 24.3  & 125.6 & 0.00796 \\
            \hline
            4  & 80  & 8.5   & 9.7   & 111.1 & 0.00900 \\
            \hline
            5  & 90  & -2.8  & -2.1  & 99.5  & 0.01005 \\
            \hline
            6  & 100 & -13.2 & -11.6 & 89.1  & 0.01122 \\
            \hline
            7  & 110 & -20.1 & -19.8 & 82.0  & 0.01220 \\
            \hline
            8  & 120 & -27.5 & -27.5 & 74.4  & 0.01344 \\
            \hline
        \end{tabular}    
    \end{table}

\end{flushleft}


\begin{flushleft}
    9. Расчет результатов косвенных измерений (таблицы, примеры расчетов). \\

    \begin{center}
        \textbf{Таблица 2. Зависимость углового коэффициента графика Vц(1/p) от температуры газа.} \\
    \end{center}
    \begin{table}[h]
        \centering
        \begin{tabular}{|c|c|c|}
            \hline
            № п.п. & $ t $, $^\circ C $ & $ K $, $ \mu_k $ \\
            \hline
            1  & 21.7  & 7983 ± 42 \\
            \hline
            2  & 30.5  & 8235 ± 38 \\
            \hline
            3  & 40.7  & 8472 ± 45 \\
            \hline
            4  & 49.8  & 8716 ± 50 \\
            \hline
            5  & 60.0  & 8940 ± 55 \\
            \hline
        \end{tabular}
    \end{table}
    Температура абсолютного нуля: \\
    $t_* = -\frac{C}{A} = -298.9 \ ^\circ C $ \\
    Угловой коэффициент A и свободный член C: \\
    $ A = 24.95 \ \mu k/^\circ C $ \\
    $ C = 7457.73 \ \mu k $ \\

    \newpage
    \begin{center}
        \textbf{Таблица 3. Зависимость давления газа от температуры при разных значениях объема.} \\
    \end{center}
    \begin{table}[h]
        \centering
        \begin{tabular}{|c|c|c|c|c|c|c|c|c|}
        \hline
        $ V_c $, мл & 50 & 60 & 70 & 80 & 90 & 100 & 110 & 120 \\
        \hline
        $ t, \ ^\circ C $ & \multicolumn{8}{c|}{$p$, кПа} \\
        \hline
        21.7 & 156.3 & 131.5 & 114.0 & 99.6  & 89.0  & 80.2  & 73.2  & 66.7 \\
        \hline
        30.5 & 156.4 & 132.9 & 115.2 & 102.9 & 91.4  & 82.4  & 74.7  & 68.4 \\
        \hline
        40.7 & 158.9 & 136.4 & 117.1 & 104.1 & 95.1  & 84.0  & 76.6  & 69.6 \\
        \hline
        49.8 & 165.6 & 139.0 & 120.2 & 106.6& 95.9 & 87.1 & 79.2 & 71.7 \\
        \hline
        60.0 & 170.9 & 146.0 & 125.6 & 111.1& 99.5 & 89.1 & 82.0 & 74.4 \\
        \hline
        $ 1/V_c$, мл$^{-1}$ & 0.020 & 0.0167 & 0.0143 & 0.0125 & 0.0111 & 0.0100 & 0.0091 & 0.0083 \\
        \hline
        $ t_{*}, \, ^\circ C $ & -361.80 & -332.50 & -360.69 & -336.10 & -313.05 & -320.36 & -293.84 & -319.06 \\
        \hline
        \end{tabular}
    \end{table}

    $ A' \approx -5201.15 $ \\
    $ C' \approx -263.35 $ \\
    $ E \approx 243.37 $ \\
\end{flushleft}

\begin{flushleft}
    10. Расчет погрешностей измерений (для прямых и косвенных измерений). \\
    Приборные погрешности: $ \Delta V = 1$ мл, $\Delta p = 0,1 $, кПа \\
    $ \Delta A = 1.53 \ \mu k/^\circ C $ \\
    $ \Delta C = 146.3 \ \mu k $ \\
    $ \Delta C' \approx 19.59 $ \\
    $ \Delta t_* \approx \Delta C' \approx 19.59 $ \\
\end{flushleft}

\begin{flushleft}
    11. Графики (перечень графиков, которые составляют Приложение 2). \\
    \includegraphics[height=9cm]{chart(1).png} \\
    \includegraphics[height=9cm]{chart(2).png} \\
    \includegraphics[height=11cm]{chart(3).png} \\
\end{flushleft}

\begin{flushleft}
    12. Окончательные результаты. \\
    $ A = 24.95 \ \mu k/^\circ C $ \\
    $ C = 7457.73 \ \mu k $ \\
    $t_* = -298.9 \ ^\circ C $ \\
    $\Delta t_* = 19.59  $ \\
    $ A' \approx -5201.15 $ \\
    $ C' \approx -263.35 $ \\
\end{flushleft}


\newpage


\begin{flushleft}
    13. Выводы и анализ результатов работы. \\
    \hspace{0.63cm}На основе проведенных измерений зависимости давления $p$ от рабочего
    объема $V_c$ при различных температурах $t$ (от 20°C до 60°C) были построены графики
    $V_c(1/p)$ для разных значений температуры и $p(t)$ для различных объемов.
    Установлено, что зависимость $V_c(1/p)$ при постоянной температуре является линейной,
    что подтверждает приближенное соответствие воздуха уравнению состояния идеального
    газа Менделеева-Клапейрона.
\end{flushleft}

\begin{flushleft}
    14. Дополнительные задания. \\
\end{flushleft}

\begin{flushleft}
    15. Выполнение дополнительных заданий. \\
\end{flushleft}

\begin{flushleft}
    16. Замечания преподавателя (исправления, вызванные замечаниями преподавателя, также помещают в этот пункт). \\
\end{flushleft}

\end{document}